\section{Additional Diagnostics on Sentiment Disagreement}
\label{app:disagreement_diagnostics}

This appendix provides additional diagnostics on the disagreement between the
GPT-based sentiment measure and the Japanese LMMD-based sentiment measure.
The goal is to distinguish economically meaningful disagreement from mechanical
differences arising from scale, centering, and observations close to neutrality.

\subsection{Motivation}
\label{app:disagreement_motivation}

Although the two sentiment measures are positively correlated in the cross-section,
they need not capture identical aspects of disclosure tone. The GPT-based measure
is designed to reflect overall document-level sentiment, whereas the LMMD-based
measure aggregates token-level positive and negative word usage. As a result,
disagreement may arise for at least three reasons. First, the two measures may
differ in scale and dispersion. Second, they may differ in how neutrality is
centered around zero. Third, they may differ substantively in what they detect. For example, one
method may capture discourse-level negativity while the other is influenced by
boilerplate managerial language or mixed segment reporting.

To separate these channels, we construct both continuous and discrete measures of
disagreement and report additional distributional and case-based diagnostics.

\subsection{Continuous Disagreement Measures}
\label{app:disagreement_continuous}

Let $\text{GPT}_{i,t}$ denote the document-level GPT sentiment score for filing
$i$ at time $t$, and let $\text{LMMD}_{i,t}$ denote the LMMD net sentiment score.
Because the two measures are expressed on different raw scales, our primary
continuous disagreement measure is based on standardized scores:
\begin{align}
z^{\text{GPT}}_{i,t} &= \frac{\text{GPT}_{i,t} - \overline{\text{GPT}}}{s_{\text{GPT}}}, \\
z^{\text{LMMD}}_{i,t} &= \frac{\text{LMMD}_{i,t} - \overline{\text{LMMD}}}{s_{\text{LMMD}}},
\end{align}
where $\overline{\text{GPT}}$ and $\overline{\text{LMMD}}$ are sample means and
$s_{\text{GPT}}$ and $s_{\text{LMMD}}$ are sample standard deviations.

We then define signed and absolute disagreement as:
\begin{align}
\text{Disagreement}^{z}_{i,t} &= z^{\text{LMMD}}_{i,t} - z^{\text{GPT}}_{i,t}, \\
\left| \text{Disagreement}^{z}_{i,t} \right| &= 
\left| z^{\text{LMMD}}_{i,t} - z^{\text{GPT}}_{i,t} \right|.
\end{align}

The absolute standardized disagreement measure,
$\left| \text{Disagreement}^{z}_{i,t} \right|$, is our preferred continuous
proxy for the magnitude of disagreement because it is invariant to differences
in the raw ranges of the two sentiment metrics.

For transparency, we also retain the raw-scale analogues:
\begin{align}
\text{Disagreement}^{\text{raw}}_{i,t} &= \text{LMMD}_{i,t} - \text{GPT}_{i,t}, \\
\left| \text{Disagreement}^{\text{raw}}_{i,t} \right| &=
\left| \text{LMMD}_{i,t} - \text{GPT}_{i,t} \right|.
\end{align}
However, these raw-scale measures are interpreted cautiously because the two
underlying sentiment scores are not naturally commensurate.

\subsection{Directional Disagreement and Neutral Bands}
\label{app:disagreement_directional}

A natural discrete notion of disagreement is whether the two methods assign
opposite signs to the same filing. A raw sign-based conflict indicator is:
\begin{equation}
\text{Conflict}^{0}_{i,t} =
\mathbb{1}\!\left[
\left( \text{GPT}_{i,t} > 0 \ \wedge \ \text{LMMD}_{i,t} < 0 \right)
\ \vee \
\left( \text{GPT}_{i,t} < 0 \ \wedge \ \text{LMMD}_{i,t} > 0 \right)
\right].
\end{equation}

However, raw sign conflict may overstate substantive disagreement if many
observations lie close to zero. To address this issue, we define method-specific
neutral bands and construct a thresholded polarity-conflict indicator.

For the GPT score, we classify filings as:
\begin{equation}
\text{Class}^{\text{GPT}}_{i,t} =
\begin{cases}
\text{positive} & \text{if } \text{GPT}_{i,t} > \varepsilon_{\text{GPT}}, \\
\text{negative} & \text{if } \text{GPT}_{i,t} < -\varepsilon_{\text{GPT}}, \\
\text{neutral}  & \text{if } |\text{GPT}_{i,t}| \le \varepsilon_{\text{GPT}},
\end{cases}
\end{equation}
where $\varepsilon_{\text{GPT}}$ is set to a fixed threshold.

For the LMMD score, we analogously define:
\begin{equation}
\text{Class}^{\text{LMMD}}_{i,t} =
\begin{cases}
\text{positive} & \text{if } \text{LMMD}_{i,t} > \varepsilon_{\text{LMMD}}, \\
\text{negative} & \text{if } \text{LMMD}_{i,t} < -\varepsilon_{\text{LMMD}}, \\
\text{neutral}  & \text{if } |\text{LMMD}_{i,t}| \le \varepsilon_{\text{LMMD}},
\end{cases}
\end{equation}
where $\varepsilon_{\text{LMMD}}$ is calibrated from the empirical distribution
of the LMMD score.

The thresholded polarity-conflict indicator is then:
\begin{equation}
\text{Conflict}^{\ast}_{i,t} =
\mathbb{1}\!\left[
\left( \text{Class}^{\text{GPT}}_{i,t} = \text{positive}
\ \wedge \
\text{Class}^{\text{LMMD}}_{i,t} = \text{negative} \right)
\ \vee \
\left( \text{Class}^{\text{GPT}}_{i,t} = \text{negative}
\ \wedge \
\text{Class}^{\text{LMMD}}_{i,t} = \text{positive} \right)
\right].
\end{equation}

This thresholded measure is more conservative than the raw sign-conflict
indicator and is intended to capture stronger directional disagreement.

\subsection{Distributional Diagnostics}
\label{app:disagreement_distributional}

To assess whether zero is a meaningful empirical center for each sentiment
measure, we report the mean, median, selected quantiles, and the shares of
negative, neutral, and positive observations for both GPT and LMMD sentiment.
These diagnostics help determine whether apparent disagreement is driven by
substantive interpretive differences or by differences in centering near zero.

In addition, we compare the raw sign-conflict rate to the thresholded
polarity-conflict rate. A large decline from $\text{Conflict}^{0}_{i,t}$ to
$\text{Conflict}^{\ast}_{i,t}$ indicates that much of the apparent disagreement
comes from observations near neutrality rather than strong directional
contradiction.

\subsection{Robustness to GPT Validity Filters}
\label{app:disagreement_validity}

Because a small number of GPT observations may reflect incomplete or degenerate
chunk-level outputs, we also report disagreement diagnostics after restricting
the sample to filings with valid GPT extraction. This robustness check helps
ensure that our disagreement measures are not driven by failed or empty GPT
processing.

Specifically, we recompute the correlation between GPT and LMMD sentiment, the
mean and median of absolute standardized disagreement, and both the raw and
thresholded conflict rates in the valid-GPT subsample.

\subsection{Illustrative Filing-Level Examples}
\label{app:disagreement_examples}

These examples are drawn from filings with high values of the thresholded
polarity-conflict indicator and are intended to illustrate recurring sources of
disagreement between GPT-based and lexicon-based sentiment measures. Taken
together, they suggest that disagreement often arises when document-level
narrative tone, strategic or boilerplate managerial language, and segment-level
heterogeneity pull the two measures in different directions.

\paragraph{Example 1: JFE Holdings (E01264, 5411.T; filing date: 2020-06-19).}
Scores: GPT = -0.67; LMMD net = 0.0074; 
$|\mathrm{Disagreement}^{z}| = 3.34$.

This filing provides the clearest example of cross-method disagreement. GPT
assigns sharply negative sentiment, consistent with the document’s repeated
emphasis on severe business conditions, profit deterioration, impairment losses,
and the emergence of net losses. By contrast, the LMMD score remains slightly
positive, plausibly because the filing also contains substantial strategic and
managerial language associated with growth, strengthening, improvement, and
corporate value creation, as well as mixed segment-level discussion.
Table~\ref{tab:jfe_disagreement_phrases} illustrates these competing textual
signals.

\begin{table}[htbp]
\centering
\caption{Illustrative phrases from JFE Holdings MD\&A (filing date 2020-06-19)}
\label{tab:jfe_disagreement_phrases}
\begin{threeparttable}
\begin{tabular}{p{4.0cm} p{5.0cm} p{6.5cm}}
\toprule
\textbf{Japanese phrase} & \textbf{English gloss} & \textbf{Interpretation} \\
\midrule
\multicolumn{3}{l}{\textit{Negative narrative / adverse performance language}} \\
\jp{極めて厳しい経営環境} & extremely severe business environment & Strong document-level negative framing. \\
\jp{大幅に悪化} & deteriorated substantially & Marked worsening in performance. \\
\jp{赤字となりました} & became loss-making & Direct signal of negative earnings. \\
\jp{減損損失を計上} & recorded impairment losses & Negative accounting event. \\
\jp{需要の低迷} & demand slump & Macro and industry deterioration. \\
\jp{目標の達成は困難} & target achievement is difficult & Negative forward-looking assessment. \\
\addlinespace
\multicolumn{3}{l}{\textit{Positive / strategic language likely supporting the lexicon score}} \\
\jp{持続的な成長} & sustainable growth & Positive strategic wording. \\
\jp{企業価値の向上} & corporate value enhancement & Standard positive managerial language. \\
\jp{成長戦略の推進} & promotion of growth strategy & Strategic positive phrasing. \\
\jp{収益拡大} & profit expansion & Positive business vocabulary. \\
\jp{強化} & strengthening & Common positive management term. \\
\jp{改善} & improvement & Positive token. \\
\bottomrule
\end{tabular}
\begin{tablenotes}[flushleft]
\footnotesize
\item Notes: The table reports selected phrases from the MD\&A of JFE Holdings
that help explain the disagreement between GPT-based and lexicon-based sentiment.
The upper block highlights phrases contributing to a strongly negative
document-level interpretation, while the lower block reports positive or
strategic language that may mechanically support a mildly positive lexicon score.
\end{tablenotes}
\end{threeparttable}
\end{table}

\paragraph{Example 2: Asahi Group Holdings (E00394, 2502.T; filing date 2020-03-26).}
Scores: GPT = -0.14; LMMD net = 0.0155;
$|\mathrm{Disagreement}^{z}| = 3.33$.

This filing provides a different type of disagreement from Example~1. GPT assigns
mildly negative sentiment, which is consistent with the document's repeated
discussion of declines in consolidated revenue and profit, adverse weather
conditions, foreign-exchange headwinds, competitive pressure, and weaker
performance in selected businesses. By contrast, the LMMD score remains
positive, plausibly because the filing is heavily interwoven with strategic and
managerial language concerning brand strengthening, premiumization, efficiency
improvements, growth initiatives, and value creation. In this case, the
disagreement appears to arise not from a sharply distressed narrative, but from
the coexistence of modest performance deterioration with an otherwise confident
and forward-looking managerial tone. Table~\ref{tab:asahi_disagreement_phrases} illustrates these competing signals.

\begin{table}[htbp]
\centering
\caption{Illustrative phrases from Asahi Group Holdings MD\&A (filing date 2020-03-26)}
\label{tab:asahi_disagreement_phrases}
\begin{threeparttable}
\begin{tabular}{p{4.0cm} p{5.0cm} p{6.5cm}}
\toprule
\textbf{Japanese phrase} & \textbf{English gloss} & \textbf{Interpretation} \\
\midrule
\multicolumn{3}{l}{\textit{Negative / weaker-performance language}} \\
\jp{減収} & revenue decline & Weaker top-line performance. \\
\jp{減益} & profit decline & Negative earnings signal. \\
\jp{競争激化} & intensifying competition & Pressure. \\
\addlinespace
\multicolumn{3}{l}{\textit{Positive / strategic language likely supporting the lexicon score}} \\
\jp{価値創造} & value creation & Strong positive managerial phrasing. \\
\jp{収益構造改革} & profit structure reform & Improvement language. \\
\jp{ブランド価値の向上} & enhancement of brand value & Positive managerial language. \\
\bottomrule
\end{tabular}
\begin{tablenotes}[flushleft]
\footnotesize
\item Notes: The table reports selected phrases from the MD\&A of Asahi Group Holdings
that help explain the disagreement between GPT-based and lexicon-based sentiment.
The upper block highlights weaker performance and adverse operating conditions,
while the lower block reports strategic and brand-oriented language that may
mechanically support a positive lexicon score.
\end{tablenotes}
\end{threeparttable}
\end{table}

\paragraph{Example 3: Alps Alpine Co., Ltd. (E01793, 6770.T; filing date: 2022-06-23).}
Scores: GPT = 0.50; LMMD net = -0.0089;
$|\mathrm{Disagreement}^{z}| = 2.89$.

This filing illustrates a third and distinct source of disagreement. Unlike
Example~1, which reflects a strongly negative document-level narrative despite
some positive strategic wording, and Example~2, where modest deterioration
coexists with forward-looking managerial language, the present case is better
characterized as a recovery narrative layered over substantial operating
frictions and segment-level weakness. GPT assigns positive sentiment, plausibly
because the filing emphasizes an overall rebound in economic conditions,
recovery in selected end markets, higher consolidated sales and profits, and a
return to net income. By contrast, the LMMD score is slightly negative, likely
because the same text repeatedly invokes supply-chain disruption, semiconductor
shortages, unstable component supply, cost inflation, sluggish display-related
performance, and an operating loss in the car information equipment segment.
Table~\ref{tab:example3_disagreement_phrases} illustrates these opposing
signals.

\begin{table}[htbp]
\centering
\caption{Illustrative phrases from Alps Alpine Co., Ltd. MD\&A (filing date 2022-06-23)}
\label{tab:example3_disagreement_phrases}
\begin{threeparttable}
\begin{tabular}{p{4.6cm} p{4.2cm} p{5.8cm}}
\toprule
\textbf{Japanese phrase} & \textbf{English gloss} & \textbf{Interpretation} \\
\midrule
\multicolumn{3}{l}{\textit{Positive narrative / improved-results language likely supporting GPT positivity}} \\
\jp{景気は回復傾向となりました} & economic recovery trend & Positive macro recovery framing. \\
\jp{好転の兆しを見せました} & signs of improvement & Positive turnaround language. \\
\jp{売上高が増加しました} & sales increased & Direct positive performance signal. \\
\addlinespace
\multicolumn{3}{l}{\textit{Negative operating / risk language likely weighing on the LMMD score}} \\
\jp{物流費や部材等の高騰} & soaring costs & Cost pressure. \\
\jp{サプライチェーンの混乱は深刻さを増しており} & worsening supply-chain disruption & Strong operating headwind. \\
\jp{営業損失は45億円} & 4.5 bn yen operating loss  & Direct negative earnings signal. \\
\bottomrule
\end{tabular}
\begin{tablenotes}[flushleft]
\footnotesize
\item Notes: The table reports selected phrases from the MD\&A of Alps Alpine Co., Ltd. that help explain the disagreement between GPT-based and lexicon-based sentiment. The upper block highlights recovery framing and improved results that may support a positive document-level GPT assessment. The lower block highlights cost pressure, disruption, and losses that may weigh on the LMMD score.
\end{tablenotes}
\end{threeparttable}
\end{table}

\subsection{Relation to the Main Analysis}
\label{app:disagreement_relation}

The diagnostics in this appendix inform the interpretation of the main
cross-sectional tests in two ways. First, they justify the use of
$\left| \text{Disagreement}^{z}_{i,t} \right|$ as the main continuous measure of
cross-method divergence. Second, they motivate the use of
$\text{Conflict}^{\ast}_{i,t}$ as a conservative indicator of strong directional
disagreement. Together, these measures allow us to distinguish between mild
differences around neutrality and more economically meaningful disagreement in
the interpretation of Japanese MD\&A disclosures.